% =====================================================================
%  1bat-ud6-5.tex  —  SESSIÓ 5: Probabilitat condicionada: quan una cosa depèn d'una altra
%  (sense preàmbul: s'inclou des de main.tex)
% =====================================================================
\horatitol{Sessió 5 — Probabilitat condicionada: quan una cosa depèn d'una altra}%
{Una probabilitat pot canviar quan ja sabem alguna cosa: introduïm la probabilitat condicionada i la distinció entre successos dependents i independents.}

\subsection{Idea clau}

Arribem al concepte més delicat de la unitat (per això la programació la situa al
tercer trimestre): la \textbf{probabilitat condicionada}. La idea és intuïtiva: saber
una cosa pot \textbf{canviar} la probabilitat d'una altra. Per exemple, la probabilitat
que un usuari trobi feina canvia si sabem que ha completat un programa de formació.

\begin{idea}
\textbf{Probabilitat condicionada}\\
$P(A\mid B)$ es llegeix «probabilitat d'$A$ \textbf{sabent que} ja ha passat $B$». La
fórmula:
\[
P(A\mid B)=\frac{P(A\cap B)}{P(B)}.
\]
\textbf{Dependents o independents:}
\begin{itemize}[leftmargin=1.4em, itemsep=2pt]
  \item \textbf{Dependents:} saber $B$ \emph{canvia} la probabilitat d'$A$ (és a dir,
        $P(A\mid B)\neq P(A)$).
  \item \textbf{Independents:} saber $B$ \emph{no} canvia res ($P(A\mid B)=P(A)$). En
        aquest cas, $P(A\cap B)=P(A)\per P(B)$.
\end{itemize}
\textbf{El més difícil no és el càlcul, sinó llegir bé l'enunciat:} identificar quina
informació es dóna com a \emph{ja coneguda}.
\end{idea}

\subsection{Exemples}

\begin{exemple}[la probabilitat que canvia amb la informació]
En un programa, de cada $100$ joves, $70$ completen la formació i $30$ l'abandonen.
\textbf{Dels que completen la formació}, $56$ troben feina; \textbf{dels que
abandonen}, només $9$. La probabilitat de trobar feina canvia segons el que sabem:
\[
P(\text{feina}\mid\text{completa})=\frac{56}{70}=0,8=80\%,\qquad
P(\text{feina}\mid\text{abandona})=\frac{9}{30}=0,3=30\%.
\]
Saber si ha completat la formació \textbf{canvia molt} la probabilitat de trobar feina
($80\%$ contra $30\%$): els dos successos són \textbf{dependents}.
\end{exemple}

\begin{exemple}[dependents contra independents]
\textbf{Dependents:} treure dues cartes d'una baralla \emph{sense tornar la primera}.
La probabilitat de la segona depèn de quina va sortir primer (queden menys cartes).

\textbf{Independents:} tirar una moneda i després un dau. Que surti cara \textbf{no}
afecta gens el que farà el dau: $P(\text{6}\mid\text{cara})=P(\text{6})=\frac{1}{6}$.
Quan dos successos no s'afecten, són independents.
\end{exemple}

\subsection{Exercicis resolts}

\begin{exresolt}[calcular una condicionada amb la fórmula]
Se sap que $P(A\cap B)=0,2$ i $P(B)=0,5$. Calcula $P(A\mid B)$.
\solucio
Apliquem la fórmula:
\[
P(A\mid B)=\frac{P(A\cap B)}{P(B)}=\frac{0,2}{0,5}=0,4=40\%.
\]
Sabent que ha passat $B$, la probabilitat d'$A$ és del $40\%$.
\resposta{$P(A\mid B)=0,4=40\%$.}
\end{exresolt}

\begin{exresolt}[llegir l'enunciat: quina és la condició]
En un casal, $P(\text{fa esport})=0,5$ i $P(\text{fa esport i música})=0,15$. Quina és
la probabilitat que una persona faci esport \textbf{sabent que} fa música, si
$P(\text{música})=0,3$?
\solucio
Ens demanen $P(\text{esport}\mid\text{música})$. La condició («sabent que») és
\textbf{música}, de manera que va al denominador:
\[
P(\text{esport}\mid\text{música})=\frac{P(\text{esport i música})}{P(\text{música})}
=\frac{0,15}{0,3}=0,5=50\%.
\]
En aquest cas $P(\text{esport}\mid\text{música})=0,5=P(\text{esport})$: saber que fa
música \emph{no} canvia la probabilitat de fer esport, de manera que són
\textbf{independents}.
\resposta{$0,5$; a més, esport i música són independents.}
\end{exresolt}

\subsection{Exercicis per fer}

\begin{exercicis}
\begin{exlist}
  \item Amb $P(A\cap B)=0,12$ i $P(B)=0,4$, calcula $P(A\mid B)$.
  \item Amb $P(A\cap B)=0,3$ i $P(A)=0,6$, calcula $P(B\mid A)$. (Compte: ara la
        condició és $A$.)
  \item En un grup, $P(\text{treballa})=0,6$ i $P(\text{treballa i estudia})=0,18$.
        Calcula $P(\text{estudia}\mid\text{treballa})$.
  \item Digues si cada parella de successos és \textbf{dependent} o \textbf{independent}
        i justifica-ho:
    \begin{enumerate}[label=\alph*), leftmargin=1.6em, topsep=2pt]
      \item tirar dos daus diferents;
      \item treure dues boles d'una bossa sense tornar la primera;
      \item que plogui i que una persona porti paraigua.
    \end{enumerate}
  \item \textbf{(Interpretació)} $P(A\mid B)$ i $P(B\mid A)$ són sempre iguals? Posa un
        exemple de l'àmbit social on \emph{no} ho siguin (pensa: «probabilitat de ser
        usuari del servei sabent que és menor» contra «probabilitat de ser menor sabent
        que és usuari»).
  \item \textbf{(Raonament)} Per què, en un problema de probabilitat condicionada, la
        part més difícil sol ser \textbf{llegir l'enunciat} i no fer el càlcul? Quina
        paraula de l'enunciat sol indicar la condició?
\end{exlist}
\end{exercicis}

% ---------------------------------------------------------------------
\fullsolucions{Sessió 5}
\begin{solucions}
\begin{sollist}
  \item $P(A\mid B)=\dfrac{0,12}{0,4}=0,3=30\%$.
  \item $P(B\mid A)=\dfrac{P(A\cap B)}{P(A)}=\dfrac{0,3}{0,6}=0,5=50\%$.
  \item $P(\text{estudia}\mid\text{treballa})=\dfrac{0,18}{0,6}=0,3=30\%$.
  \item (a) \textbf{Independents}: el resultat d'un dau no afecta l'altre. \quad
        (b) \textbf{Dependents}: en no tornar la primera bola, canvien els casos
        possibles per a la segona. \quad (c) \textbf{Dependents}: si plou, és més
        probable que porti paraigua; saber una cosa canvia la probabilitat de l'altra.
  \item \textbf{No}, en general $P(A\mid B)\neq P(B\mid A)$. Exemple: la probabilitat de
        ser usuari del servei \emph{sabent que} és menor pot ser baixa (pocs menors el
        fan servir), mentre que la probabilitat de ser menor \emph{sabent que} és
        usuari pot ser alta (si el servei és sobretot per a joves). Són preguntes
        diferents.
  \item Resposta oberta. Idea: el càlcul és una simple divisió; la dificultat és
        identificar \textbf{quina informació es dóna com a coneguda} (la condició, que
        va al denominador) i quina és la que es pregunta. La paraula que sol indicar la
        condició és «\textbf{sabent que}» (o «si», «donat que», «entre els que\dots»).
\end{sollist}
\end{solucions}
